\originalchapter{恒星轨道、导引中心与局部旋转测量}

本章对应 8.902 第4讲末、第5讲与第6讲开头. 已知共同势后, 每颗恒星的轨道由初值与守恒量决定. 球对称势的运动可降到一个平面; 轴对称盘则通常只保留轴向角动量. 这两种对称性不同, 正是圆轨道、径向摆动和竖直摆动能够分开近似的前提.

\originalsection{守恒量减少自由变量, 不自动关闭轨道}
定常势中单位质量能量 $E=v^2/2+\Phi$ 守恒, 因为 $\dd E/\dd t=\partial_t\Phi=0$. 球对称势的加速度沿径向, 所以单位质量角动量 $\boldsymbol\ell=\boldsymbol r\times\boldsymbol v$ 也守恒. 非零 $\boldsymbol\ell$ 把运动限制在垂直于它的平面. 取该平面的极坐标 $(r,\psi)$,
\[
\boldsymbol v=\dot r\hat{\boldsymbol e}_r+r\dot\psi\hat{\boldsymbol e}_\psi,
\quad \ell=r^2\dot\psi,
\quad E=\tfrac12\dot r^2+\Phi_{\rm eff}(r),
\quad\Phi_{\rm eff}=\Phi+\frac{\ell^2}{2r^2}.
\]
径向转折点满足 $E=\Phi_{\rm eff}$, 径向加速度 $\ddot r=-\Phi_{\rm eff}'$. 角变量仍通过 $\dot\psi=\ell/r^2$ 重建. 当 $\ell=0$ 时是径向运动, 不宜再用角变量参数化轨道.

只保留轨道平面时有 $(r,\psi,\dot r,\dot\psi)$ 四个变量, 再固定 $E,\ell$ 后可进一步降维. 原课把“有守恒量”直接说成四维相空间, 容易混淆消去方向、固定守恒量与整个系综的维数. 在一般轴对称势中, $E,\ell_z$ 也不保证所有轨道可积; 不能凭两个积分写出普遍第三积分.

\originalsection{Binet 方程与闭合条件}
令 $u=1/r$, 由于 $\dd/\dd t=\ell u^2\dd/\dd\psi$, 有 $\dot r=-\ell u'$ 和 $\ddot r=-\ell^2u^2u''$. 代入 $\ddot r-r\dot\psi^2=-\Phi'(r)$,
\[
u''+u=\frac{\Phi'(1/u)}{\ell^2u^2}.
\]
这不是“去掉 $r=0$ 奇点”的保证, 而是在非零角动量区用角度取代时间. Kepler 势 $\Phi=-GM/r$ 给常数右端 $GM/\ell^2$, 所以
\[
u=\frac{GM}{\ell^2}[1+e\cos(\psi-\psi_0)],\qquad
r=\frac{p}{1+e\cos(\psi-\psi_0)},\quad p=\ell^2/(GM).
\]
对 $0\le e<1$, 是闭合椭圆, $E=-GM/(2a)$ 与 $p=a(1-e^2)$. 一般势中的径向周期与角向周期未必可公度, 轨道常为进动玫瑰形. 近圆轨道的一径向周期积累角度 $2\pi\Omega/\kappa$; 只有合适频率比时轨道闭合.

课程用 $\Phi=-GM/r-2G^2M^2/(c^2r^2)$ 演示进动. 该玩具中心势给
\[
u''+\left(1-\frac{4G^2M^2}{\ell^2c^2}\right)u=\frac{GM}{\ell^2},
\qquad \Delta\psi\simeq\frac{4\pi GM}{pc^2}.
\]
它可以说明频率改变会产生进动, 但不等于 Schwarzschild 的完整弱场轨道方程. 后者给 $6\pi GM/(pc^2)$, 含速度依赖的相对论动力学, 见已有 GR 卷. 本卷不重复 GR 推导, 也不把上述玩具模型声称为精确水星进动公式. 原课的振子式中位移与速度导数混排, 此处以明确的 $u''+k^2u$ 恢复其示范目的.

\originalsection{轴对称势的二维有效运动}
取光滑定常势 $\Phi(R,z)$, 不依赖方位角 $\varphi$, 则 $\ell_z=R^2\dot\varphi$ 守恒. 柱坐标单位向量的时间导数给
\[
\ddot R-R\dot\varphi^2=-\partial_R\Phi,
\quad R\ddot\varphi+2\dot R\dot\varphi=0,
\quad\ddot z=-\partial_z\Phi.
\]
定义 $\Phi_{\rm eff}=\Phi+\ell_z^2/(2R^2)$ 后,
\[
\ddot R=-\partial_R\Phi_{\rm eff},\quad
\ddot z=-\partial_z\Phi_{\rm eff},\quad
E=\tfrac12(\dot R^2+\dot z^2)+\Phi_{\rm eff}.
\]
可达区域由 $E\ge\Phi_{\rm eff}$ 限制, 方位运动由瞬时 $R$ 重建, 因而通常不是角速度严格均匀的圆周运动. 若势对 $z\mapsto-z$ 对称且光滑, 盘面圆轨道 $R=R_g,z=0$ 满足
\[
\partial_R\Phi(R_g,0)=\frac{\ell_z^2}{R_g^3},
\qquad \Omega_g^2=\frac{\partial_R\Phi(R_g,0)}{R_g}.
\]
导引中心 $R_g$ 是与给定 $\ell_z$ 相应的圆轨道半径, 不是恒星每时每刻的实际位置. 圆轨道是否稳定, 还需有效势二阶导数的符号.

\originalsection{本轮近似由有效势 Hessian 得到}
令 $x=R-R_g$, 在光滑对称盘面附近展开. 一阶项由圆轨道条件消失, 混合项由反射对称消失, 因而
\[
\Phi_{\rm eff}\simeq\Phi_{\rm eff}(R_g,0)+\tfrac12\kappa^2x^2+\tfrac12\nu_z^2z^2,
\quad
\kappa^2=\Phi_{RR}+3\Omega_g^2,\quad \nu_z^2=\Phi_{zz}.
\]
$\kappa$ 为径向本轮频率, $\nu_z$ 为竖直频率, 后者加下标避免与辐射频率混淆. 比较 $\Omega^2(R)=\Phi_R(R,0)/R$, 得
\[
\kappa^2=R\frac{\dd\Omega^2}{\dd R}+4\Omega^2
=\frac1{R^3}\frac{\dd}{\dd R}(R^4\Omega^2).
\]
两频率平方都正才有两个稳定小振动. 取 $x=X\cos(\kappa t+\alpha)$, $z=Z\cos(\nu_zt+\beta)$, 又由角动量守恒
\[
\dot\varphi\simeq\Omega_g(1-2x/R_g),\quad
\varphi\simeq\Omega_gt+\varphi_0-
\frac{2\Omega_gX}{\kappa R_g}\sin(\kappa t+\alpha).
\]
随导引中心旋转的局部切向坐标 $y=R_g(\varphi-\Omega_gt-\varphi_0)$, 故 $y=-2\Omega_gX\sin(\kappa t+\alpha)/\kappa$. 平面内运动是逆向的本轮椭圆, 长短轴比由 $2\Omega_g/\kappa$ 决定.

\begin{example}
若 $v_c\propto R^q$, 则 $\Omega\propto R^{q-1}$, $\kappa^2=2(q+1)\Omega^2$. Kepler 区 $q=-1/2$ 给 $\kappa=\Omega$; 平坦曲线 $q=0$ 给 $\kappa=\sqrt2\Omega$; 刚体旋转 $q=1$ 给 $\kappa=2\Omega$. 若 $q<-1$, 径向频率平方负, 对应圆轨道不稳定, 不应仍画成稳定小椭圆.
\end{example}

\begin{center}
\begin{tikzpicture}[>=Latex]
\draw[->](-1.6,0)--(1.6,0)node[right]{$x/X$};
\draw[->](0,-1.7)--(0,1.7)node[above]{$y/X$};
\draw[structurecolor,thick,domain=0:360,samples=120,variable=\t]plot({cos(\t)},{-sqrt(2)*sin(\t)});
\draw[->,structurecolor,thick](1,0)--(0.98,-0.25);
\node[below right] at(0,0){导引中心};
\end{tikzpicture}
\end{center}
图为平坦曲线的小振幅极限, 箭头取从 $x=X$ 开始的运动方向. 无限薄 Mestel 盘的势在 $z=0$ 有 $|z|$ 尖角, $\Phi_{zz}$ 不是普通连续函数, 不满足这个光滑竖直 Taylor 近似. 真实有限厚盘或其他光滑模型才能直接使用有限 $\nu_z$.

\originalsection{Oort 常数需要区分速度与角速度}
在太阳附近作 $d\ll R_0$ 的局部展开, 假设扣除太阳特有运动后恒星随近圆轨道运动. 定义
\[
A=-\tfrac12R_0\Omega'_0=\tfrac12(v_c/R_0-v_c'),
\qquad B=-\Omega_0-\tfrac12R_0\Omega'_0.
\]
因此 $A-B=\Omega_0$, $\kappa_0^2=-4B(A-B)$. $A,B$ 都是速度除长度的量纲; 原笔记的 $A$ 缺少 $R_0$ 时不满足量纲.

若银经 $l$ 从朝银心方向起算, 银纬为 $b$, 正视线速度指远离太阳, 则局部一阶表达式为
\[
v_{\rm los}\simeq dA\sin(2l)\cos^2b,\qquad
v_l\simeq d[A\cos(2l)+B]\cos b.
\]
相应沿银经的大圆角速度 $\mu_l^*\equiv\dot l\cos b=v_l/d$; 故在盘面 $b=0$, $\mu_l^*=A\cos(2l)+B$, 而不是再乘一个距离. 原课用 $\mu$ 表示带 $d$ 的式子, 若解释为角自行会量纲错误; 此处分别记 $v_l$ 与 $\mu_l^*$. 实际拟合还要考虑距离误差、恒星速度弥散、非轴对称流和非圆轨道, 局部公式不是整个银河的通用速度场.

\originalsection{光度—速度幂律中的额外假设}
维里尺度 $v^2\sim GM/R$, 质光比 $\Upsilon=M/L$, 平均面亮度尺度 $\Sigma_L\sim L/R^2$ 合并给
\[
v^4\sim G^2\Upsilon^2\Sigma_LL,\qquad
L\sim\frac{v^4}{G^2\Upsilon^2\Sigma_L}.
\]
若同类星系 $\Upsilon,\Sigma_L$ 及结构系数近似相同, 才得到 $L\propto v^4$. 课程以此解释螺旋星系 Tully--Fisher 与椭圆星系 Faber--Jackson 关系的尺度动机. 前者比较旋转速度, 后者比较弥散, 不能互换测量定义. 经验斜率、散布和波段依赖需由样本标定, 不是维里定理单独固定. 一旦标定可由速度估计光度, 再与通量比较求距离, 这把本章动力学与第1章观测连接起来.
